%% Abstract -- rewritten 2026-09-29 to the Copernicus length guideline (<= 250 words).
\begin{abstract}
A DOAS fit reports a residual, chi-square, uncertainty and convergence flag, read as a measure of retrieval
quality. We present Augur, an open retrieval that poses the fit in the extinction domain as a
separable nonlinear least-squares problem, solves it by variable projection and is configured by instrument and
campaign profiles. Its calibration references enter as data, so they can be perturbed, and the objective mapped, without
reprocessing a campaign. We use it on one field case, a
thermal-dissociation broadband cavity-enhanced absorption spectrometer \citep{Min_2016,Nam_2022} with two heated and one
ambient-temperature cavity, to measure two gaps between fit report and data.

The first is error the fit cannot see. For the alkyl nitrate product, leave-one-out perturbation of the zero-air,
reflectivity and etalon references gives, in the best-calibrated week, a structural uncertainty \fieldnum{1.28} to \fieldnum{2.29}
times the reported per-scan sigma, and a total \fieldnum{1.63} to \fieldnum{2.50} times, without enlarging the residual. The second is information the pipeline holds and
discards. The ambient-temperature channel's wavelength shift is not identifiable from ambient spectra, yet its fits
report convergence. A time-base error introduced by the pipeline, a reflectivity reference nine
hours out of alignment, displaced the product by \fieldnum{0.100}\,ppb while changing the residual and the reported uncertainty
by two parts in a thousand. A multiplicative scale error, evident only against a co-located reference analyser, is invisible to both.
We propose that each record carry its termination state, distance to the nearest calibration and reference
provenance, and release a synthetic benchmark for other implementations.
\end{abstract}
